%% Dans ce fichier, les �ventuels accents sont cod�s
%% en ISO-8859-15 (aka latin9)
%% La ligne suivante donne trois exemples 
%% e accent aigu : �, e accent grave : �, e-dans-l'o : �
%%%%%ENTREES%%%%
\begin{tabarg}
\multicolumn{3}{l}{\titretab {Entr\'ees}}\\
\hline


$X$ & Vecteur, \texttt{float[$N_{level}$]} & Signal lidar\\
$X_{gate}$ & Coeff., \texttt{string[$N_{level}$]} & Portes du lidar (\$Arg1.bornes) ($m$)\\
$Z_{lid}$ & Vecteur, \texttt{float[$N_{level}$]} & Altitudes r�elles des portes du lidar ($m$)\\
$N_{level}$ & Coeff., \texttt{float} & Nombre de portes (\$Arg1.array\_dim)\\
$Z_{p}$& Vecteur, \texttt{float[$N_p$]} & Altitude des niveaux du profil de $T$ et $P$ (m)\\
$T$ & Vecteur, \texttt{float[$N_p$]} & Profil de temp�rature ($K$ ou $Celsius$)\\
$P$ & Vecteur, \texttt{float[$N_p$]} & Profil de pression($hPa$)\\
$N_p$ & Coeff., \texttt{Int} & Nombre de niveaux du profil de $T$ et $P$ ( \$Arg6.array\_dim)\\
$\lambda$ & Coeff., \texttt{float} & Longueur d'onde du signal lidar ($nm$)\\
\end{tabarg}

\medskip
%%%%%SORTIES%%%%
\begin{tabarg}
\multicolumn{3}{l}{\titretab {Sorties}}\\
\hline

$\alpha_{mol}$ & Vecteur, \texttt{float[$N_{level}$]} & Profil du coefficient d'extinction mol�culaire ($m^{-1}$)\\
$\beta_{mol}$ & Vecteur, \texttt{float[$N_{level}$]} &  Profil du coefficient de r�tro-diffusion mol�culaire ($m^{-1}$)\\
$T_{rayleigh}$ & Vecteur, \texttt{float[$N_{level}$]} & Profil du coefficient de Rayleigh \\
$Pr^2_{mol}$ & Vecteur, \texttt{float[$N_{level}$]} & Signal mol�culaire r�tro-diffus� ($Photons.m^2$)\\

\end{tabarg}
